\section*{Capítulo \ref{ch:g11:enzymes-and-phenotype} --- As enzimas e o fenótipo}

\begin{solution}{exo:g11:enzymes-and-phenotype:1}
\omterm{def:g11:enzymes-and-phenotype:phenotype}{Genótipo}: o conjunto de \omterm{def:g10:universal-dna:gene}{alelos} carregados. \omterm{def:g11:enzymes-and-phenotype:phenotype}{Fenótipo}: o conjunto das
características observáveis, nas escalas molecular, celular e do
organismo.
\end{solution}

\begin{solution}{exo:g11:enzymes-and-phenotype:2}
O bolso da \omterm{def:g11:enzymes-and-phenotype:enzyme}{enzima} em que o \omterm{def:g11:enzymes-and-phenotype:enzyme}{substrato} se liga e é transformado. A
especificidade: o sítio se ajusta a um \omterm{def:g11:enzymes-and-phenotype:enzyme}{substrato} (ou a uma família de
\omterm{def:g11:enzymes-and-phenotype:enzyme}{substratos}) e catalisa uma reação.
\end{solution}

\begin{solution}{exo:g11:enzymes-and-phenotype:3}
Perto de \qty{37}{\celsius}; a \qty{45}{\celsius} resta cerca de um terço.
\end{solution}

\begin{solution}{exo:g11:enzymes-and-phenotype:4}
Molecular: tirosinase ausente ou inativa, nenhuma melanina fabricada.
Celular: \omterm{def:g10:cells-common-unit:cell}{células} do pigmento presentes, mas vazias de melanina. Do
organismo: cabelos brancos, pele e olhos muito claros, sensibilidade à
luz e às queimaduras de sol.
\end{solution}

\begin{solution}{exo:g11:enzymes-and-phenotype:5}
O dobramento da pepsina só funciona em ácido forte (pH 2); o intestino
é neutro a levemente alcalino, onde a curva mostra atividade zero.
\end{solution}

\begin{solution}{exo:g11:enzymes-and-phenotype:6}
Um \omterm{def:g10:universal-dna:gene}{alelo} funcional produz metade da \omterm{def:g11:enzymes-and-phenotype:enzyme}{enzima} normal, mas uma molécula de
\omterm{def:g11:enzymes-and-phenotype:enzyme}{enzima} trabalha milhares de vezes por segundo: metade da quantidade
ainda converte tirosina suficiente para fabricar pigmento normal. O
\omterm{def:g11:enzymes-and-phenotype:phenotype}{fenótipo} está saturado; só dois \omterm{def:g10:universal-dna:gene}{alelos} que falham aparecem.
\end{solution}

\begin{solution}{exo:g11:enzymes-and-phenotype:7}
O albinismo deles vem de bloqueios em dois \omterm{def:g10:universal-dna:gene}{genes} diferentes (duas
\omterm{def:g11:enzymes-and-phenotype:enzyme}{enzimas} diferentes da via, ou do transporte da melanina). A criança
herda um \omterm{def:g10:universal-dna:gene}{alelo} funcional de cada \omterm{def:g10:universal-dna:gene}{gene} do genitor cujo bloqueio está no
outro, de modo que as duas \omterm{def:g11:enzymes-and-phenotype:enzyme}{enzimas} funcionam e o pigmento é fabricado.
\end{solution}

\begin{solution}{exo:g11:enzymes-and-phenotype:8}
A fenilalanina se acumula; a tirosina (e o que deriva dela) está
faltando. Como a melanina é fabricada a partir da tirosina, a falta
dela reduz o pigmento: cabelos e pele claros são comuns na
fenilcetonúria não tratada.
\end{solution}

\begin{solution}{exo:g11:enzymes-and-phenotype:9}
Onde o leite é um alimento básico, digeri-lo fornece energia e cálcio
sem adoecer, de modo que o \omterm{def:g10:universal-dna:gene}{alelo} persistente é uma vantagem. Onde os
adultos não bebem leite, o \omterm{def:g10:universal-dna:gene}{alelo} não muda nada; ele é ``mais saudável''
apenas no ambiente que fornece o \omterm{def:g11:enzymes-and-phenotype:enzyme}{substrato}.
\end{solution}

\begin{solution}{exo:g11:enzymes-and-phenotype:10}
Para a tirosinase: molecular, o \omterm{def:g11:enzymes-and-phenotype:enzyme}{substrato} não encaixa mais, nenhuma
melanina; celular, \omterm{def:g10:cells-common-unit:cell}{células} do pigmento vazias; do organismo, albinismo.
Para a lactase: lactose não digerida; \omterm{def:g10:cells-common-unit:cell}{células} do intestino privadas do
açúcar e o intestino grosso fermentando-o; intolerância.
\end{solution}

\begin{solution}{exo:g11:enzymes-and-phenotype:11}
As extremidades dele ficam mais quentes que \qty{35}{\celsius}, a
\omterm{def:g11:enzymes-and-phenotype:enzyme}{enzima} quase não trabalha, e o filhote fica claro por inteiro. Ao ar
livre no inverno as orelhas e as patas esfriam, a \omterm{def:g11:enzymes-and-phenotype:enzyme}{enzima} fica ativa, e
a pelagem que cresce ali nasce escura.
\end{solution}

\begin{solution}{exo:g11:enzymes-and-phenotype:12}
Molecular: metade da hemoglobina é a forma falciforme. Celular: com
oxigênio normal as \omterm{def:g10:cells-common-unit:cell}{células} continuam redondas; com oxigênio muito baixo
algumas falcizam. Do organismo: saudável na vida comum, em risco em
altitude ou no mergulho --- e protegido contra a malária. O ambiente (o
oxigênio, a malária) decide se o \omterm{def:g11:enzymes-and-phenotype:phenotype}{fenótipo} do portador é um fardo ou uma
vantagem.
\end{solution}

\begin{solution}{exo:g11:enzymes-and-phenotype:13}
A \qty{41}{\celsius} as \omterm{def:g11:enzymes-and-phenotype:enzyme}{enzimas} ainda conservam a maior parte da
atividade delas (a curva está perto do topo); a \qty{43}{\celsius} a
atividade cai a pique à medida que as \omterm{def:g11:gene-expression:protein}{proteínas} começam a se desdobrar
e, se o desdobramento avança, ele é irreversível: as \omterm{def:g11:enzymes-and-phenotype:enzyme}{enzimas} são
destruídas, e não apenas desaceleradas.
\end{solution}

\begin{solution}{exo:g11:enzymes-and-phenotype:14}
Mais \omterm{def:g10:cells-common-unit:organelle}{mitocôndrias} e mais capilares por fibra muscular; um coração maior
com volume sistólico maior; mais hemoglobina por litro de sangue;
reservas maiores de \omterm{def:g10:exercise-and-energy:glycogen}{glicogênio}. Cada uma é um \omterm{def:g11:enzymes-and-phenotype:phenotype}{fenótipo} produzido pelo
mesmo \omterm{def:g11:enzymes-and-phenotype:phenotype}{genótipo} num ambiente diferente: as características
``genéticas'' são a faixa que o \omterm{def:g11:enzymes-and-phenotype:phenotype}{genótipo} permite, e não um valor fixo.
\end{solution}

\begin{solution}{exo:g11:enzymes-and-phenotype:15}
O \omterm{def:g11:enzymes-and-phenotype:phenotype}{genótipo} fixa que \omterm{def:g11:enzymes-and-phenotype:enzyme}{enzima} é fabricada e como ela se comporta: nenhuma
\omterm{def:g11:enzymes-and-phenotype:enzyme}{enzima} 1 (fenilcetonúria), nenhuma tirosinase (albinismo), uma
tirosinase sensível ao calor (siamês). O ambiente só consegue mudar o
resultado dentro do que a \omterm{def:g11:enzymes-and-phenotype:enzyme}{enzima} permite: uma dieta retira o \omterm{def:g11:enzymes-and-phenotype:enzyme}{substrato}
e salva o cérebro, mas não fabrica a \omterm{def:g11:enzymes-and-phenotype:enzyme}{enzima} que falta; o frio faz a
\omterm{def:g11:enzymes-and-phenotype:enzyme}{enzima} siamesa trabalhar, mas nenhuma temperatura faz a do albino
trabalhar; e nada no ambiente dá pigmento ao albino. O \omterm{def:g11:enzymes-and-phenotype:phenotype}{genótipo} fixa a
faixa; o ambiente escolhe dentro dela.
\end{solution}

\begin{solution}{pb:g11:enzymes-and-phenotype:1}
\textbf{1.} Tronco \qty{38}{\celsius}: cerca de 2\%, claro. Pernas
\qty{36}{\celsius}: cerca de 20\%, claro. Patas \qty{33}{\celsius}:
80\%, escuro. Orelhas e focinho \qty{31}{\celsius}: mais de 90\%, escuro.

\textbf{2.} Corpo e pernas claros, patas, orelhas e focinho (e cauda)
escuros: o padrão siamês.

\textbf{3.} 100\% em toda parte: um gato uniformemente pigmentado.

\textbf{4.} No útero toda parte do filhote está a \qty{38}{\celsius}, a
\omterm{def:g11:enzymes-and-phenotype:enzyme}{enzima} fica inativa em toda parte, e nenhum pigmento é depositado antes
do nascimento; as pontas escurecem à medida que as extremidades esfriam
depois do nascimento.

\textbf{5.} Por volta de \qty{35.5}{\celsius}, onde a atividade cruza
os 30\%: a parte descendente e íngreme da curva entre 33 e
\qty{37}{\celsius} fixa a fronteira.

\textbf{6.} Escura: a \qty{30}{\celsius} a \omterm{def:g11:enzymes-and-phenotype:enzyme}{enzima} está plenamente ativa
nos pelos em crescimento.

\textbf{7.} Clara: a \qty{38}{\celsius} a \omterm{def:g11:enzymes-and-phenotype:enzyme}{enzima} está inativa.

\textbf{8.} O pigmento é depositado no pelo enquanto ele cresce e não
pode ser retirado nem acrescentado depois; a cor só muda quando os
pelos caem e são substituídos por outros, crescidos na temperatura
local.

\textbf{9.} O \omterm{def:g11:enzymes-and-phenotype:phenotype}{genótipo} das \omterm{def:g10:cells-common-unit:cell}{células} do pedaço é o mesmo antes e depois,
e é o mesmo do resto do tronco; só a temperatura da \omterm{def:g11:enzymes-and-phenotype:enzyme}{enzima} mudou, e com
ela o pigmento. A ação é sobre a atividade da \omterm{def:g11:enzymes-and-phenotype:enzyme}{enzima}, um evento
molecular.

\textbf{10.} Todas as \omterm{def:g10:cells-common-unit:cell}{células} do gato descendem de um óvulo e carregam
os mesmos \omterm{def:g10:universal-dna:gene}{alelos}; e os experimentos do pedaço raspado mudam a cor de
uma região sem mudar as \omterm{def:g10:cells-common-unit:cell}{células} dela --- um \omterm{def:g10:universal-dna:gene}{alelo} não pode ser trocado
por uma bolsa de gelo.

\textbf{11.} Molecular: uma tirosinase que só funciona abaixo de cerca
de \qty{35}{\celsius}. Celular: \omterm{def:g10:cells-common-unit:cell}{células} do pigmento cheias de melanina
nas extremidades, vazias no tronco. Do organismo: um gato claro com as
pontas escuras.

\textbf{12.} O \omterm{def:g10:chemistry-of-life:families}{aminoácido} fica onde estabiliza o dobramento; a troca
mantém o dobramento firme em temperatura baixa, mas o deixa afrouxar
alguns graus acima, de modo que o \omterm{def:g11:enzymes-and-phenotype:enzyme}{sítio ativo} perde a forma antes de a
\omterm{def:g11:enzymes-and-phenotype:enzyme}{enzima} comum perder a dela.

\textbf{13.} A \omterm{def:g11:enzymes-and-phenotype:enzyme}{enzima} do \omterm{def:g10:universal-dna:gene}{alelo} comum é plenamente ativa na temperatura
do corpo; metade da quantidade normal ainda é mais que suficiente para
fabricar pigmento em toda parte. O \omterm{def:g10:universal-dna:gene}{alelo} siamês é recessivo em relação
a ele.

\textbf{14.} Molecular: nenhuma \omterm{def:g11:enzymes-and-phenotype:enzyme}{enzima} funcional em temperatura
nenhuma, contra uma que funciona condicionalmente. Celular: nenhuma
melanina em parte alguma, contra melanina nas regiões frias. Do
organismo: um gato todo branco com olhos rosados, contra o padrão de
pontas escuras. O frio não muda nada para o albino.

\textbf{15.} Que eles carregam um \omterm{def:g10:universal-dna:gene}{alelo} sensível à temperatura do mesmo
\omterm{def:g10:universal-dna:gene}{gene}; provavelmente uma substituição parecida que desestabiliza a
\omterm{def:g11:enzymes-and-phenotype:enzyme}{enzima}, surgida de modo independente em cada \omterm{def:g10:biodiversity-scales:species}{espécie}.

\textbf{16.} Vinte vezes; a \omterm{def:g11:enzymes-and-phenotype:enzyme}{enzima} 1, que converte a fenilalanina em
tirosina.

\textbf{17.} O excesso de fenilalanina interfere especificamente nas
\omterm{def:g10:cells-common-unit:cell}{células} do cérebro em desenvolvimento (no suprimento de outros
\omterm{def:g10:chemistry-of-life:families}{aminoácidos} delas e na maturação delas); as \omterm{def:g10:cells-common-unit:cell}{células} da pele o toleram.
O defeito molecular é o mesmo em toda parte; a consequência celular não é.

\textbf{18.} Um quinto de \num{1200} são \qty{240}{\micro mol/L},
abaixo de \num{360}: segura. Sem a \omterm{def:g11:enzymes-and-phenotype:enzyme}{enzima}, o nível no sangue apenas
acompanha a ingestão, de modo que o \omterm{def:g11:enzymes-and-phenotype:phenotype}{fenótipo} fica determinado pelo
fornecimento de \omterm{def:g11:enzymes-and-phenotype:enzyme}{substrato} --- o ambiente --- e não pela \omterm{def:g11:enzymes-and-phenotype:enzyme}{enzima} ausente
sozinha.

\textbf{19.} A fenilalanina é necessária para construir as \omterm{def:g11:gene-expression:protein}{proteínas} do
corpo e não pode ser sintetizada; sem nada dela, o crescimento para. A
ingestão precisa bastar para a síntese de \omterm{def:g11:gene-expression:protein}{proteínas} e nada mais, o que
exige medir e ajustar.

\textbf{20.} O \omterm{def:g11:enzymes-and-phenotype:phenotype}{genótipo} fixou a \omterm{def:g11:enzymes-and-phenotype:enzyme}{enzima} (sensível ao calor no gato,
ausente na criança); o ambiente determinou as condições dela (a
temperatura da pele; a fenilalanina da dieta); o \omterm{def:g11:enzymes-and-phenotype:phenotype}{fenótipo} é o resultado
do \omterm{def:g11:enzymes-and-phenotype:phenotype}{genótipo} expresso num dado ambiente.
\end{solution}
